\pagebreak
\section*{Part A: Entrepreneurial Reflection Statement}
\addcontentsline{toc}{section}{Part A: Entrepreneurial Reflection Statement}

My Master Project addressed a real-world and market problem faced by many international companies and multicultural teams, specifically when dealing with Japanese counterparts. Meetings that involve Japanese speakers and non-Japanese speakers usually have difficulties in communication due to language barriers. Another issue is the cost of hiring a translator, especially for daily meetings, which is quite high, and usually the translator is not from the specific domain, causing miscommunication to often occur. Therefore, a practical solution is needed to help teams understand Japanese speech in real time and to produce usable transcripts immediately after the meeting.

The potential target users for this solution are multinational companies which have operations in Japan or work with Japanese counterparts. Japanese companies with foreign employees or global teams can utilize this solution. Universities from Japan that conduct bilingual discussions, and conferences or webinars, will also find this solution useful, especially universities that have a huge number of international students. In particular, organizations that hold frequent meetings are strong targets because the recurring communication problem creates recurring demand.

The main feature that can be commercialized is a platform that uses the results from this project and then uses the model to create a real-time transcription during online meetings. This system can capture the audio from the system audio input and automatically create the transcript. Based on the user's default language, the transcript will be translated and then appear on the user's screen as a subtitle. In addition, the system can support post-meeting outputs such as a full transcript, timestamps, speaker turns, and searchable meeting notes. Based on the results from this project, two types of modes can be chosen by the user based on user requirements. If they prefer \textit{accuracy-first} mode, high-accuracy transformer-based models are the right option. However, if \textit{low-latency} is important, using faster CRDNN models allows teams to select the best trade-off for their use case.



The solution can be scaled or adapted for wider adoption by packaging it as a web application and meeting plugin that integrates with common platforms such as Zoom, Google Meet, and Microsoft Teams, and also by offering an on-premise deployment option for privacy-sensitive organizations. Furthermore, the system can be adapted to different domains by supporting terminology customization and glossary injection so that technical vocabulary, product names, and proper nouns can be recognized more accurately. Over time, the platform can expand from transcription into a more complete meeting intelligence tool by adding functions such as automatic summarization, action-item extraction, and bilingual outputs for cross-border collaboration.

The unique value proposition compared to existing solutions is the focus on Japanese language support with high accuracy and low latency, which is not commonly found in general-purpose transcription services. By leveraging the findings from this project, the platform can offer tailored solutions that address the specific challenges of Japanese speech recognition, such as handling honorifics, context-dependent meanings, and various dialects. Additionally, by providing flexible options for accuracy versus speed, the platform can cater to different user needs, making it more versatile than one-size-fits-all solutions.


\section*{Part B: Draft of Conclusion and Future Work Chapter }
As the conclusion, the key findings and contributions of this research project are summarized based on the three main objectives that were set out at the beginning. First, the key requirements for constructing a speech-to-text model within the context of the Japanese language were identified. This included standardizing audio formats, chunking audio files, filtering noise, and implementing tokenization for accurate WER calculation. Second, an analysis of different speech-to-text models was conducted to determine the most effective pre-processing and training setups. The study compared traditional Kaldi GMM--HMM, hybrid CRDNN--CTC, and fine-tuned Whisper models, identifying the best configurations for each. Finally, the performance of these models was evaluated using CER, WER, and RTF metrics, revealing that fine-tuned Whisper models achieved the highest accuracy, while CRDNN--CTC provided the lowest latency. These findings contribute to the understanding of effective strategies for Japanese speech-to-text transcription and provide a foundation for future developments in this area.   


From the research project, several limitations and constraints were encountered. One major limitation was the dataset used, which primarily consisted of read speech and lacked diversity in terms of conversational styles, accents, and background noise conditions. This may have affected the generalizability of the models to real-world scenarios. Additionally, computational resources were a constraint, particularly when fine-tuning larger Whisper models, which limited the extent of experimentation with model sizes and hyperparameters. The evaluation metrics used, while standard, may not fully capture user satisfaction or practical usability in real-time applications.

possible service from this project are using the developed speech-to-text models to create a real-time transcription platform for Japanese language meetings. This platform could be offered as a SaaS solution, targeting multinational corporations, educational institutions, and conference organizers who require accurate and efficient transcription services. The service could include features such as live subtitles during meetings, post-meeting transcripts, and integration with popular video conferencing tools.

From comercial perspective, the potential market for a Japanese speech-to-text transcription service is significant, given Japan's status as a major global economy with numerous multinational corporations and a growing emphasis on remote work and virtual meetings. Organizations that frequently engage in cross-border communication would find value in a reliable transcription service to enhance understanding and collaboration. Revenue could be generated through subscription models, pay-per-use pricing, or enterprise licensing agreements. Additionally, partnerships with video conferencing platforms could expand reach and user adoption.

The business model for this service could follow a SaaS approach, offering tiered subscription plans based on usage volume, features, and support levels. A freemium model could also be implemented, providing basic transcription services for free while charging for premium features such as higher accuracy models, faster processing times, and additional integrations. A startup pathway could involve initial development and testing with pilot customers, followed by iterative improvements based on user feedback. Marketing efforts would focus on digital channels, industry events, and partnerships to build brand awareness and attract early adopters.

For future improvements, the research project could explore the incorporation of more diverse datasets that include conversational speech, various dialects, and noisy environments to enhance model robustness. Additionally, experimenting with advanced acoustic models and decoding strategies could further improve accuracy and latency. Implementing user feedback mechanisms to continuously refine the models based on real-world usage would also be beneficial. Finally, expanding the platform's capabilities to include multilingual support and additional features such as speaker identification and sentiment analysis could increase its value proposition.


\section{Part C: Self-Evaluation \& Value Proposition Checklist}
\begin{itemize}
  \item[$\checkmark$] My conclusion clearly summarizes the main contributions
  \item[$\checkmark$] I have identified a real-world use case or market problem
  \item[$\checkmark$] The business element or commercialization idea is stated clearly
  \item[$\checkmark$] Future work includes potential productization or startup directions
  \item[$\checkmark$] I reflected on possible improvements or funding opportunities
\end{itemize}

Reflection 1: The strongest commercialization potential in my project lies in developing a real-time transcription platform specifically tailored for Japanese language meetings. This platform addresses a significant market need among multinational corporations and educational institutions that frequently engage in cross-border communication. By leveraging the high accuracy and low latency of the developed speech-to-text models, the service can provide substantial value to users who require reliable transcription services. The ability to integrate with popular video conferencing tools and offer features such as live subtitles and post-meeting transcripts further enhances its appeal. The SaaS business model, with tiered subscription plans and potential partnerships, provides a scalable and sustainable revenue stream, making it an attractive opportunity for commercialization.

Reflection 2: To make my project market-ready, I want to strengthen the dataset diversity and model robustness further. Expanding the dataset to include conversational speech, various dialects, and noisy environments will enhance the models' ability to perform accurately in real-world scenarios. This is crucial for ensuring that the transcription service meets the practical needs of users across different contexts. Additionally, I aim to refine the user experience by implementing feedback mechanisms that allow for continuous improvement based on real-world usage. This will help in tailoring the service to better meet user expectations and increase satisfaction. Finally, exploring advanced acoustic models and decoding strategies will be essential to optimize both accuracy and latency, ensuring that the platform remains competitive in the market.